\documentclass[10pt,a4paper]{amsart}
\usepackage[margin=19mm]{geometry}
\usepackage{amsmath,amssymb,mathtools}
\usepackage[scr=rsfs,cal=euler]{mathalfa}
\usepackage{xfrac}
\usepackage[unicode,hidelinks,bookmarks=false]{hyperref}
\newcommand{\ie}{i.e.\ }
\newcommand{\cf}{cf.\ }
\newcommand{\resp}{respectively\ }
\newcommand{\Subsection}{\subsection}
\providecommand{\nobreakdash}{\nobreak\hbox{-}}
\setlength{\emergencystretch}{2em}
\multlinegap0pt
\usepackage{bm}
\usepackage{bbm}
\usepackage{dsfont}
\usepackage{xspace}
\usepackage{cancel}
\usepackage{mathtools}
\usepackage{amsthm}
\makeatletter
\@ifundefined{thm}{\newtheorem{thm}{Thm}}{}
\@ifundefined{prop}{\newtheorem{prop}[thm]{Proposition}}{}
\makeatother
\title[Examples of left-orderable and non-left-orderable HNN extens]{Examples of left-orderable and non-left-orderable HNN extensions}
\author{Azer Akhmedov, Cody Martin}
\date{Source: arXiv:2301.00052v3 (CC BY 4.0)}
\begin{document}
\maketitle

\noindent\emph{Source and licence.} Examples of left-orderable and non-left-orderable HNN extensions. Version arXiv:2301.00052v3 (CC BY 4.0), distributed under the Creative Commons Attribution 4.0 International licence, as recorded in the arXiv licence metadata for this exact version and shown on the abstract page https://arxiv.org/abs/2301.00052v3. Full rights notice: licenses/arxiv-2301.00052-CC-BY-4.0.txt.\par\smallskip

\noindent\textit{Editorial scope.} Counted unit: the Prop whose source label is \texttt{prop:extension} in the article (source0.tex lines 281--285, source label \texttt{prop:extension}), delivered together with the article's own complete original proof of it (source0.tex lines 287--307, one \texttt{proof} environment of 21 lines). No mathematics is written, repaired or re-derived: statement and proof are the article's own text line for line. Required context retained uncounted: none. Every definition, display and label of those lines is retained unchanged. Omitted from this package and not redistributed: 1--280, 286--286, 308--410. Nothing inside a retained excerpt is omitted or altered: every retained body is delivered exactly as the article's lines are written, so no omission receipt of any kind accompanies this package. Citations: the retained excerpts carry no citation command at all, so no bibliography entry of the article is transcribed and no reference is redistributed. Reference adaptation: none was needed and none was made; every reference inside the delivered unit resolves inside it, so no unresolved reference or citation is produced. Printed slips kept literal: no printed word, symbol, hypothesis or number of the counted statement or of its proof was altered. Layout and class: the article's own class is \texttt{amsart} with packages including amsmath, amssymb, euscript, graphicx, hyperref, verbatim; the wrapper is the lane's pinned \texttt{amsart} editorial wrapper at 10pt on A4, which restates the theorem-like environments the retained text needs and adds the article's own macro definitions that the retained text needs (none), with extra wrapper packages (bm, bbm, dsfont, xspace, cancel, mathtools). The delivered numbering is the unit's own. Assets: no retained span contains a figure environment, a graphics-inclusion command, a PDF-page inclusion, a table or a TikZ picture; no image file is copied into this package, the package declares no asset dependency and no image file is redistributed with it. Licence: version arXiv:2301.00052v3 of this article is distributed under the Creative Commons Attribution 4.0 International licence, as recorded in the arXiv licence metadata for this exact version and shown on the abstract page https://arxiv.org/abs/2301.00052v3. A full rights notice is delivered at licenses/arxiv-2301.00052-CC-BY-4.0.txt. Characters: every retained excerpt is pure ASCII, so the delivered engine, XeLaTeX with its default Latin Modern text font, reports no missing character.
  \begin{prop}\label{prop:extension}  Let  \( \mathfrak{g}, \mathfrak{h} \subseteq \mathfrak{u}_n \) be finite-dimensional nilpotent Lie algebras embedded in \(\overline{ \mathfrak{u}} \), and let \( \phi : \mathfrak{g} \to \mathfrak{h} \) be a Lie algebra isomorphism. Then there exists an automorphism \( \Phi \in \mathrm{Aut}(\overline{\mathfrak{u}}) \) such that
\[
\Phi|_{\mathfrak{g}} = \phi.
\]
\end{prop}

\begin{proof}
Since \( \mathfrak{g} \) and \( \mathfrak{h} \) are finite-dimensional and isomorphic, choose disjoint finite index intervals \( I = [i_1, i_2] \subset \mathbb{Z} \) and \( J = [j_1, j_2] \subset \mathbb{Z} \), both of length \( n \), such that $|x-y| > n$ for all $x\in I, y\in J$ and 
\[
\mathfrak{g} \hookrightarrow \mathfrak{u}_I,\quad \mathfrak{h} \hookrightarrow \mathfrak{u}_J,
\]
where \( \mathfrak{u}_I \subset \overline{\mathfrak{u}} \) denotes the subalgebra of matrices supported on rows and columns in \( I \), and similarly for \( \mathfrak{u}_J \).

\medskip 

Define \( \iota_{\mathfrak{g}} : \mathfrak{g} \hookrightarrow \mathfrak{u}_I \) and \( \iota_{\mathfrak{h}} : \mathfrak{h} \hookrightarrow \mathfrak{u}_J \) as the chosen embeddings and let $\psi : \iota_{\mathfrak{g}}(\mathfrak{g})\to \iota_{\mathfrak{h}}(\mathfrak{h})$ defined as $\psi =                \iota_{\mathfrak{h}}\circ \phi \circ \iota_{\mathfrak{g}}^{-1}$. It suffices to prove that $\psi $ can be lifted to an isomorphism of \(\overline{ \mathfrak{u}} \). 

\medskip 

 Since the index sets $I$ and $J$ have length $n$ and are at a distance at least $n$, the isomorphism $\psi $ can be lifted to $\psi ':\mathfrak{u}_I\to \theta ^{k,l}(\mathfrak{u}_J)$ for some $0\leq k, l\leq n-1$.  

\medskip 

 Let also $\mathcal{B}_0, \mathcal{B}_1, \mathcal{B}$ be bases of $\iota_{\mathfrak{g}}(\mathfrak{g}), \mathfrak{u}_I, \overline{\mathfrak{u}}$ respectively, such that $\mathcal{B}_0\subseteq \mathcal{B}_1\subset \mathcal{B}$. Then we extend the isomorphism $\psi '$ (which is already defined in $\mathcal{B}_1$) inductively along the basis $\mathcal{B}$ and obtain an isomorphism $\Psi : \overline{\mathfrak{u}}\to \overline{\mathfrak{u}}$. This extension is straightforward if we use the following standard fact: For all $k\geq 1$, let $\mathcal{N}_k$ denote the Lie algebra of $k\times k$ real strictly upper-triangular matrices. $\mathcal{N}_k$ has a standard embedding in $\mathcal{N}_{k+1}$ as the sub-algebra of upper-left block of size $k\times k$.  Then any Lie algebra automorphism $\mathcal{N}_k\to \mathcal{N}_k$ lifts to an automorphism $\mathcal{N}_{k+1}\to \mathcal{N}_{k+1}$. Indeed, this fact also suggests a specific natural order to extend the basis $\mathcal{B}_1$ to $\mathcal{B}$.


\end{proof}

\end{document}
